On considère l’équation différentielle (E) : $y' + y = 1-e^{-x}$ où $y$ est une fonction de la variable réelle $x$, définie et dérivable sur l’ensemble $\R$ des nombres réels, et $y'$ est la fonction dérivée de $y$.

\begin{enumerate}
	\item Donner et résoudre l'équation homogène (E$_0$)\sautp

	\Lignes{2}
	\item A l’aide du logiciel Géogébra, déterminer les nombres réels $a$ et $b$ pour que la fonction $p$ définie sur $\R$ par : $p(x) = a + b~x~e^{-x}$ soit une solution particulière de l'équation différentielle (E), et donner l’expression de $p$. 
		
		{\it On prendra des curseurs allant de $0$ à $40$ par pas de $2$}\sautm
		
		On trouve : $a=$\makebox[0.1\linewidth]{\dotfill} \hfill $b=$\makebox[0.1\linewidth]{\dotfill} \hfill $p(x)=$\makebox[0.3\linewidth]{\dotfill}\sauts
		
		\appelprof
	
	\item En déduire la forme générale des solutions de (E).\sautp

	\Lignes{1}
	\item Déterminer, à l'aide de géogebra et d'un curseur bien choisi, la solution particulière $f$ de (E) telle que $f'(0)=0$.\sautp

	\Lignes{1}
	
	\appelprof
\end{enumerate}